\gddchapter{Interview sequence}


Below is the guideline for the semi-open guided interview. 
The interview starts with easy warm-up questions and progresses towards harder questions as it goes along. 

The answers were translated from polish back to english.

\section{Introduction}
Remind the participant that this is not a test, rather a joint exploration.

\section{Icebreakers}

\begin{QandA}
   \item Now that you’ve tried both versions, how are you feeling? Was there anything that immediately surprised you?
         \begin{answered}
               okay so um i think it was easier to play with like clicking buttons than speaking but it's because i'm not comfortable speaking i'm not used to it and i think that
               It's a pretty creative way to play games.
               I haven't played any games like this before, so it was something new for me.
               And I think it gives you the illusion of choice, because you don't know what you can do and what you cannot do.
               Yeah, I think it was great.
            \end{answered}

   \item If you had to describe the difference between these two games to a person who hasn't played either, what would you say?
         \begin{answered}
            Okay, so I would say that the second version with clicking is more simple and can be much faster because you don't have to think what you can do and what you cannot.
            You have simple choices that you can see.
            And with talking you have to use your imagination with what you can do and what you want to do.
            So I would say its second version is more limited to something.
            Yeah, I think I would explain it that way.
            Gabriel: From the moment that you used your voice to navigate, what were you thinking as you decided what to say to the game?
            I didn't know how to pick the best words because I thought to myself, will it understand me?
            So I was trying to be as clear as possible and as official as I can be.
         \end{answered}
\end{QandA}




\section{Grand tour questions}
(remember to pause here 3-5 seconds)

\begin{QandA}


   \item From the moment you first used your voice to navigate, what were you thinking as you decided what to say to the system?

         \begin{answered}
            I didn't know how to pick the best words because I thought to myself, will it understand me?
            So I was trying to be as clear as possible and as official as I can be..
         \end{answered}

    \item How was it like to use the voice input in the game?
         \begin{answered}
            It was weird.
            As I said, it's because I don't use my voice to anything else than communicate with others, with other people.
            I do not use anything that requires voice guidance.
            So I think It was weird for me, but it will be easier for people that use voice comments to communicate via Alexa or Google Assistant You don't talk to chat GPT for instance, do you?
         \end{answered}


\end{QandA}




\section{Core comparative questions}

\begin{QandA}
   \item In which version did you feel more 'inside' the story world?
         \begin{answered}
            I think in the second one, because I could look at the screen and don't have to think too much of what to pick.
            and how to create proper sentence and I think it's easier for me to pick options that are already available and not to give comments on my own because I
            I think if I had an idea of what I would be playing, then I would be more comfortable with talking, but because it was like told me on the spot, I couldn't prepare myself and I think the second option was more comfortable for me.

           
         \end{answered}

   \item How did the 'effort' of thinking of what to say naturally compare to the 'effort' of knowing which keys to press?

         \begin{answered}
            I think that I read much faster than I speak.
            I think I was doing much much better in the second option because I was just sometimes I didn't finish reading and I could pick an option for finishing reading so it was easier for me and I think I saved some time before because of that and I did better progress in the second time
         \end{answered}

    \item Did the freedom to say 'anything' make you feel more in control, or was it more overwhelming than having a fixed list of keys?

         \begin{answered}
            I think it was overwhelming, but it's because English is not my first language.
            I think it will be easier when I would
            be able to speak Polish and because I'm stressed because we're in public.
            But I think if I was in like more comfortable place for me then it would be easier for me and it wouldn't be as stressful as it was.
         \end{answered}
    \item How would you compare the overall experience of using the two game versions? Which one would you choose and why?

         \begin{answered}
            Wiktoria: On the first thought, I would choose the second one.
            But when I think about it more, I think the first option would be more challenging for me.
            And if I could test it, I don't know, in my home, in some more quiet environment, that it would be a good way to practice speaking to the machines, I guess.
            Gabriel: Yeah, yeah.
            Gabriel: I understand.
            I think it's a good way to be more comfortable with speaking.
            If you have any problems with communicating with others and you are worried of talking to other people, I think it's a good exercise for that.
            Because you need to think and properly... How to say that?
            Exonerate?
            I think I would say that.
            Your thoughts.


            
         \end{answered}
\end{QandA}






\section{Probing}
In this part of the interview look into the sticky moments noted above and ask follow-up questions probing for more detail.
Include lexical reflection (eg. you used the word "clunky" when explaining navigation. Could you tell me more about what was happening
in that moment?)

Most of the probing in this interview happened as a follow up to the question "In which version did you feel more 'inside' the story world?"

\begin{QandA}
   \item In the beginning you said that it is kind of like a chatbot. What do you mean by that?
   
      \begin{answered}
             When I was in the high school, there was this chatbot game when you picked some options and then the game, I mean the chatbot, the chatbot will show you the plot.
            You can just choose something and it will
            walk you through their plot.
            I don't know how to explain.
            I can't even remember what the game is called.
            But it was something like that.
            So I've played something similar and it was Childhood Mechanics.
            I think it was on Nintendo.
         
      \end{answered}


        
\end{QandA}


\section{Final reflection}

\begin{QandA}
   \item After reflecting on everything today, is there anything else you’d like to add about using your voice to play this game that we haven't talked about?

         \begin{answered}
            I think it's something new and it's cool to see something new.
            It's a pretty chill game to play and it's cool that you can interact with every room that you enter and you can take a look around and see what's around you.
            Yeah, I think it would be even cooler if you could do more things and if you could try to do more stuff.
            But overall, I think it's something new and something that can be revolutionary.
          
         \end{answered}
\end{QandA}

\section{Closing}
Thank the participant for his time and tell them again how much value his input and time gives me.






